\documentclass[11pt,a4paper]{article}

% ------------------------------------------------------------
% Packages
% ------------------------------------------------------------
\usepackage[a4paper,margin=22mm]{geometry}
\usepackage{graphicx}
\usepackage{booktabs}
\usepackage{array}
\usepackage{longtable}
\usepackage{siunitx}
\usepackage{amsmath,amssymb}
\usepackage{float}
\usepackage{xcolor}
\usepackage{listings}
\usepackage{caption}
\usepackage{enumitem}
\usepackage[hidelinks]{hyperref}
\usepackage{fancyhdr}
\usepackage{microtype}

\sisetup{detect-all,per-mode=symbol}
\setlength{\parindent}{0pt}
\setlength{\parskip}{0.55em}
\setlength{\headheight}{14pt}
\setlist[itemize]{noitemsep,topsep=0.3em}

\pagestyle{fancy}
\fancyhf{}
\lhead{EEE3096S Practical 1B}
\rhead{MCHTHA046 \& PLLVEL003}
\cfoot{\thepage}

% ------------------------------------------------------------
% Listings setup
% ------------------------------------------------------------
\lstdefinelanguage{ARMAsm}{
  morekeywords={LDR,STR,MOVS,MOV,LSRS,ANDS,TST,BEQ,BNE,B,BL,BX,PUSH,POP,SUBS,NOP,.syntax,.thumb,.cpu,.fpu,.global,.type,.equ,.section,.size},
  sensitive=true,
  morecomment=[l]{@}
}

\lstset{
  language=ARMAsm,
  basicstyle=\ttfamily\small,
  keywordstyle=\bfseries,
  commentstyle=\itshape,
  numbers=left,
  numberstyle=\tiny,
  stepnumber=1,
  numbersep=8pt,
  frame=single,
  breaklines=true,
  columns=fullflexible,
  showstringspaces=false
}

% ------------------------------------------------------------
% Helper commands
% ------------------------------------------------------------
\newcommand{\todoanswer}[1]{%
  \par\noindent\fcolorbox{black}{gray!5}{%
    \parbox{0.95\linewidth}{\textbf{TO FILL:} #1\vspace{1.5em}}%
  }\par
}

\newcommand{\figureplaceholder}[3]{%
\begin{figure}[H]
  \centering
  \IfFileExists{#1}{%
    \includegraphics[width=0.92\linewidth]{#1}%
  }{%
    \fbox{\parbox[c][55mm][c]{0.88\linewidth}{\centering
      \textbf{Insert figure here}\\[0.6em]
      Expected file: \texttt{\detokenize{#1}}\\[0.6em]
      #2
    }}%
  }
  \caption{#3}
\end{figure}
}

\newcommand{\codefile}[2]{%
  \IfFileExists{#1}{\lstinputlisting[caption={#2}]{#1}}{%
    \fbox{\parbox[c][35mm][c]{0.94\linewidth}{\centering
      Copy the final source file to \texttt{\detokenize{#1}} before compiling the submission.\\
      #2
    }}%
  }%
}

% ------------------------------------------------------------
% Document
% ------------------------------------------------------------
\begin{document}

\begin{titlepage}
\centering
\vspace*{20mm}
{\Large University of Cape Town\\Department of Electrical Engineering\par}
\vspace{18mm}
{\huge\bfseries EEE3096S Practical 1B\par}
\vspace{3mm}
{\Large Execution Timing and Cycle-Accurate Assembly on the UCT Dev Board\par}
\vspace{20mm}
\begin{tabular}{ll}
\textbf{Student 1:} & Thapelo Mocheko (MCHTHA046)\\
\textbf{Student 2:} & Velan Pillay (PLLVEL003)\\
\textbf{Year:} & 2026\\
\end{tabular}
\vfill
\textit{Report filename required by the practical:}\\
\texttt{prac1B\_MCHTHA046\_PLLVEL003.pdf}
\vspace{15mm}
\end{titlepage}

\tableofcontents
\clearpage

% ============================================================
\section{Golden Measure and PC Timing -- Task 1}
% Target length from practical: approximately 2 pages.

\subsection{Method and predicted execution time}
Describe the PC implementation of the golden integer square-root function using double-precision arithmetic and the standard-library square-root function. State the instruction-count estimate and the PC CPU clock used before running the timing experiment.

\todoanswer{CPU model and clock frequency; estimated instruction count; predicted time per call in ns; brief reasoning.}

\subsection{Timing runs}
The practical requires two timing runs with different repetition counts and a comparison of the resulting time per call.

\begin{table}[H]
\centering
\caption{Task 1 timing runs.}
\begin{tabular}{lrrrr}
\toprule
Run & Repetitions & Total measured time & Time/call (ns) & Notes\\
\midrule
1 & & & & \\
2 & & & & \\
\bottomrule
\end{tabular}
\end{table}

\todoanswer{Mean time per call and spread across the two runs. Define clearly how ``spread'' was calculated.}

\subsection{Golden outputs}
Complete the full table for the ten required inputs.

\begin{table}[H]
\centering
\caption{Golden integer square-root outputs.}
\begin{tabular}{r r l}
\toprule
Input $x$ & Golden output $\lfloor\sqrt{x}\rfloor$ & Verification/status\\
\midrule
0 & & \\
1 & & \\
15 & & \\
16 & & \\
4095 & & \\
65535 & & \\
123456789 & & \\
987654321 & & \\
4294836225 & & \\
4294967295 & & \\
\bottomrule
\end{tabular}
\end{table}

\subsection{Hand check}
Write one check in full by proving both inequalities
\[
 y^2 \le x
 \qquad\text{and}\qquad
 (y+1)^2 > x.
\]

\todoanswer{Choose one required input and show the complete arithmetic.}

\subsection{Task 1 question}
\todoanswer{State the golden output for 987654321, the measured time per call in ns, the repetition count used, and explain why timing one isolated call is unreliable.}

% ============================================================
\section{Board Implementation and Timer Measurement -- Task 2}
% Target length from practical: approximately 3 pages plus figures.

\subsection{Clock path and timer configuration}
Name every clock-tree stage from HSI to the selected 16-bit timer, including divider values and the timer prescaler.

\todoanswer{Timer used; HSI frequency; SYSCLK/AHB/APB path; APB timer-clock rule if applicable; prescaler value; resulting counter frequency and counter tick period. Cite the exact RM0091 sections used.}

\subsection{Predicted timing}
\todoanswer{Predicted instruction count, predicted counter span, and predicted execution time for one integer square-root call. Show reasoning.}

\subsection{Measured single-call timing}
\begin{align}
\Delta N &= \text{counter span} = \boxed{\phantom{00000}},\\
f_{\mathrm{counter}} &= \boxed{\phantom{00000000}\ \mathrm{Hz}},\\
\Delta t &= \frac{\Delta N}{f_{\mathrm{counter}}}
          = \boxed{\phantom{000.000}\ \mu\mathrm{s}}.
\end{align}

\todoanswer{Insert the actual counter values before/after the call and show the complete divider-to-microseconds calculation.}

\figureplaceholder{figures/task2_single_call.png}
{Full uncropped scope capture. PC13 pulse, channel information, V/div, time/div, both cursors and delta readout visible.}
{Task 2 single-call execution pulse on PC13. The body text must state the same measured pulse width as the cursor readout.}

\subsection{Counter wrap-around}
For a 16-bit free-running counter, one wrap occurs after $2^{16}=65536$ counts. Show the long-run arithmetic explicitly rather than assuming unsigned subtraction without explanation.

\todoanswer{Long-run repetition count; start and end counter values; number of wraps; total count span; total time; mean time per call.}

\subsection{Golden-value verification}
\todoanswer{State pass/fail against all ten Task 1 golden values and include the evidence produced by the firmware.}

\subsection{Task 2 question}
\todoanswer{State the scope pulse width in us and the recorded counter span. Show the full calculation from divider values to pulse width. State whether the system clock matches the 8 MHz HSI and explain how this was confirmed.}

% ============================================================
\section{Optimisation Flags -- Task 3}
% Target length from practical: approximately 2 pages plus figures.

\subsection{Prediction}
\todoanswer{Predicted ratio between the -O0 and -O2 execution times, written down before building -O2, with brief reasoning.}

\subsection{Optimisation results}
\begin{table}[H]
\centering
\caption{Effect of compiler optimisation on execution time and text size.}
\begin{tabular}{lrrrr}
\toprule
Flag & Text size (bytes) & Counter time/call ($\mu$s) & Scope width ($\mu$s) & Ratio to -O0\\
\midrule
-O0 & & & & 1.000\\
-O1 & & & & \\
-O2 & & & & \\
-Os & & & & \\
\bottomrule
\end{tabular}
\end{table}

\begin{table}[H]
\centering
\caption{Predicted and measured optimisation ratios.}
\begin{tabular}{lll}
\toprule
Comparison & Predicted ratio & Measured ratio\\
\midrule
-O0 / -O2 & & \\
\bottomrule
\end{tabular}
\end{table}

\figureplaceholder{figures/task3_O0.png}{Full uncropped -O0 scope capture with cursors.}{Task 3 execution pulse at \texttt{-O0}.}
\figureplaceholder{figures/task3_O1.png}{Full uncropped -O1 scope capture with cursors.}{Task 3 execution pulse at \texttt{-O1}.}
\figureplaceholder{figures/task3_O2.png}{Full uncropped -O2 scope capture with cursors.}{Task 3 execution pulse at \texttt{-O2}.}
\figureplaceholder{figures/task3_Os.png}{Full uncropped -Os scope capture with cursors.}{Task 3 execution pulse at \texttt{-Os}.}

\subsection{Disassembly evidence}
Name one optimisation transformation visible in the \texttt{-O2} disassembly and identify the source lines on which it acts.

\todoanswer{Name of transformation, source-line numbers, and a short explanation of what changed. Paste a concise disassembly excerpt below.}

\begin{lstlisting}[language={},caption={Task 3 disassembly excerpt at -O2.}]
% Paste only the relevant disassembly excerpt here.
\end{lstlisting}

\subsection{Execution-time/code-size trade-off}
\todoanswer{Explain the measured trade-off between execution time and text size across -O0, -O1, -O2 and -Os.}

\subsection{Task 3 question}
\todoanswer{State measured -O0 and -O2 times in us; measured ratio; predicted ratio; optimisation transformation; and the source lines affected.}

% ============================================================
\section{Cycle-Counted Phase Delay -- Task 4}
% Target length from practical: approximately 2 pages plus figures.

\subsection{Required phase delay}
The input frequency is \SI{1}{\kilo\hertz}, so the period is
\begin{equation}
T=\frac{1}{f}=\frac{1}{1000}=\SI{1}{\milli\second}.
\end{equation}
A \ang{45} phase shift is one eighth of a cycle, therefore
\begin{equation}
\Delta t=\frac{45^\circ}{360^\circ}T
       =\frac{1}{8}(\SI{1}{\milli\second})
       =\boxed{\SI{125}{\micro\second}}.
\end{equation}

With the system clock at \SI{8}{\mega\hertz},
\begin{equation}
T_{\mathrm{cycle}}=\frac{1}{\SI{8}{\mega\hertz}}
=\boxed{\SI{125}{\nano\second}}.
\end{equation}
Hence the target loop period is
\begin{equation}
N_{\mathrm{cycles}}=\frac{\SI{125}{\micro\second}}{\SI{125}{\nano\second}}
=\boxed{1000\ \text{cycles}}.
\end{equation}

\subsection{Assembly implementation}
The core polling loop used for the cycle-counted design is shown below. The ADC sample is read from the ADC data register and written directly to DAC1 before the software delay.

\begin{lstlisting}[caption={Core Task 4 ADC-to-DAC loop.}]
DSP_Loop:
    LDR  R0, =ADC_DR
    LDR  R1, =DAC_DHR12R1

loop:
    LDR  R2, [R0]          @ predicted 2 cycles
    STR  R2, [R1]          @ predicted 2 cycles

    MOVS R3, #248          @ 1 cycle
    NOP                    @ 1 cycle
    NOP                    @ 1 cycle

delay_loop:
    SUBS R3, R3, #1        @ 1 cycle x 248
    BNE  delay_loop        @ 247 taken + 1 not taken

    B    loop              @ predicted 3 cycles
\end{lstlisting}

\subsection{Cycle budget}
\begin{table}[H]
\centering
\caption{Task 4 cycle budget for one ADC-to-DAC update.}
\begin{tabular}{lrrrr}
\toprule
Instruction/event & Cycles/instruction & Iterations & Subtotal & Running total\\
\midrule
\texttt{LDR R2,[R0]} & 2 & 1 & 2 & 2\\
\texttt{STR R2,[R1]} & 2 & 1 & 2 & 4\\
\texttt{MOVS R3,\#248} & 1 & 1 & 1 & 5\\
\texttt{NOP} & 1 & 2 & 2 & 7\\
\texttt{SUBS} & 1 & 248 & 248 & 255\\
\texttt{BNE} taken & 3 & 247 & 741 & 996\\
\texttt{BNE} not taken & 1 & 1 & 1 & 997\\
\texttt{B loop} & 3 & 1 & 3 & \textbf{1000}\\
\bottomrule
\end{tabular}
\end{table}

Thus the predicted update period is
\begin{equation}
1000(\SI{125}{\nano\second})=\boxed{\SI{125}{\micro\second}},
\end{equation}
which corresponds to the required \ang{45} phase displacement at \SI{1}{\kilo\hertz}.

\figureplaceholder{figures/task4_phase_delay.png}
{CH1 on PB0 and CH2 on PA4. Show the full scope screen with cursors placed on corresponding points of the two waveforms.}
{Task 4 input on PB0 and DAC output on PA4. Cursor separation $\Delta t=\underline{\hspace{16mm}}\,\mu$s, corresponding to the measured phase delay.}

\subsection{Measured error and discussion}
Use the actual oscilloscope reading in the following calculation:
\begin{align}
\Delta t_{\mathrm{meas}} &= \boxed{\phantom{000.0}\ \mu\mathrm{s}},\\
\varepsilon_t &= \Delta t_{\mathrm{meas}}-\SI{125}{\micro\second},\\
\%\text{ error} &= \frac{|\varepsilon_t|}{\SI{125}{\micro\second}}\times100
                 = \boxed{\phantom{00.00}\%}.
\end{align}

Any non-zero error should be discussed using the measured result. Plausible contributors include HSI oscillator tolerance, scope cursor placement, ADC/DAC peripheral timing and the discrete sample-and-hold nature of the DAC waveform. Do not claim a contributor that is inconsistent with the observed data.

% ============================================================
\section{Analog Debugging of the LCD Enable Line -- Task 5}
% Target length from practical: approximately 3 pages plus figures.

\subsection{4-bit LCD interface}
The LCD was driven in 4-bit mode using PC15 as Enable, PC14 as RS, and PA15/PA12/PB9/PB8 as D7/D6/D5/D4. The program waits for the LCD power rail, executes the 4-bit initialisation sequence, and writes ASCII \texttt{0x41} (``A'').

\begin{lstlisting}[caption={Task 5 entry point and character write.}]
LCD_Run:
    PUSH {LR}

    BL   LCD_DelayLong
    BL   LCD_DelayLong
    BL   LCD_DelayLong
    BL   LCD_DelayLong

    BL   LCD_Init

    MOVS R0, #0x41
    BL   LCD_WriteData

hang:
    B    hang
\end{lstlisting}

The 4-bit power-up sequence used the standard pattern of three \texttt{0x3} nibbles followed by \texttt{0x2}, after which full commands such as \texttt{0x28}, \texttt{0x08}, \texttt{0x01}, \texttt{0x06}, \texttt{0x0C} and \texttt{0x80} were sent through the 4-bit write routine.

\subsection{Uncorrected Enable timing}
Before adding NOP padding, the PC15 HIGH and LOW masks were preloaded so that no unintended literal-load instruction lengthened the HIGH pulse:

\begin{lstlisting}[caption={Uncorrected Enable pulse used for the failing-case measurement.}]
LCD_Pulse:
    PUSH {R0, R1, R2, LR}
    LDR  R0, =GPIOC_BSRR
    LDR  R1, =0x00008000      @ PC15 HIGH mask
    LDR  R2, =0x80000000      @ PC15 LOW mask

    STR  R1, [R0]             @ PC15 HIGH
    @ no NOP padding here
    STR  R2, [R0]             @ PC15 LOW

    BL   LCD_DelayShort
    POP  {R0, R1, R2, PC}
\end{lstlisting}

\figureplaceholder{figures/task5_before_fix.png}
{CH1 on PC15 and CH2 on PC15\_S using the real short Enable pulse. Show that the 5 V-side signal does not provide the required valid HIGH interval.}
{Task 5 before-fix Enable waveform: PC15 (CH1) and PC15\_S (CH2).}

\subsection{Measured rise time and parasitic capacitance}
A long diagnostic square wave was used only to make the RC rising edge easy to capture. The measured time from approximately \SI{0}{\volt} to the \SI{3.5}{\volt} threshold on PC15\_S was approximately
\begin{equation}
\boxed{t_r \approx \SI{27}{\nano\second}}.
\end{equation}
\textbf{Important:} replace \SI{27}{\nano\second} everywhere in this report if the final cursor placement at exactly \SI{3.5}{\volt} gives a different value.

\figureplaceholder{figures/task5_rise_time.png}
{CH2 source selected. Place one time cursor at the beginning of the PC15\_S rise and the second at the point where PC15\_S reaches 3.5 V.}
{Measured PC15\_S rise time from 0 V to 3.5 V: $t_r=\underline{\hspace{14mm}}$ ns.}

The practical gives the capacitor-charging relation
\begin{equation}
V_C(t)=V_{CC}\left(1-e^{-t/(RC)}\right),
\end{equation}
which gives
\begin{equation}
\boxed{
C=\frac{-t_r}{R\ln\!\left(1-\dfrac{V_C(t)}{V_{CC}}\right)}
}.
\end{equation}
For this board, substitute
\[
t_r=\SI{27}{\nano\second},\qquad
V_C(t)=\SI{3.5}{\volt},\qquad
R=\boxed{\phantom{0000}\ \Omega},\qquad
V_{CC}=\boxed{\phantom{0.00}\ \mathrm{V}}.
\]
Therefore,
\begin{equation}
C=\boxed{\phantom{000.00}\ \mathrm{pF}}.
\end{equation}
The resistor value and supply voltage must be taken from the actual board/schematic and bench measurement; they are intentionally not invented here.

\subsection{NOP padding calculation}
At \SI{8}{\mega\hertz}, each CPU cycle is \SI{125}{\nano\second}. Using the measured rise time of approximately \SI{27}{\nano\second} and the practical requirement that PC15\_S remain above \SI{3.5}{\volt} for at least \SI{450}{\nano\second}, a conservative required MCU-side HIGH time is
\begin{equation}
t_{\mathrm{required}}=t_r+\SI{450}{\nano\second}
\approx \SI{27}{\nano\second}+\SI{450}{\nano\second}
=\SI{477}{\nano\second}.
\end{equation}
Hence
\begin{equation}
N_{\mathrm{NOP}}
=\left\lceil\frac{\SI{477}{\nano\second}}{\SI{125}{\nano\second}}\right\rceil
=\left\lceil3.816\right\rceil
=\boxed{4\ \mathrm{NOPs}}.
\end{equation}
The four NOPs provide a nominal padding of
\begin{equation}
4(\SI{125}{\nano\second})=\boxed{\SI{500}{\nano\second}}.
\end{equation}

\begin{lstlisting}[caption={Corrected Enable pulse with four NOPs.}]
    LDR  R1, =0x00008000
    LDR  R2, =0x80000000
    STR  R1, [R0]             @ PC15 HIGH

    NOP
    NOP
    NOP
    NOP

    STR  R2, [R0]             @ PC15 LOW
\end{lstlisting}

\subsection{Fixed timing evidence}
The final oscilloscope measurement must be made using the corrected \emph{real} LCD Enable pulse, not the long diagnostic square wave. Place the cursors on the rising and falling \SI{3.5}{\volt} crossings of PC15\_S.

\figureplaceholder{figures/task5_after_fix.png}
{CH1 on PC15 and CH2 on PC15\_S. Place the cursors at the two 3.5 V crossings of CH2. The delta must be at least 450 ns.}
{Task 5 fixed Enable waveform. PC15\_S remains above 3.5 V for $\Delta t=\underline{\hspace{14mm}}$ ns, satisfying the 450 ns requirement.}

Record the final result explicitly:
\begin{equation}
\boxed{t_{\mathrm{PC15\_S}>3.5V}=\phantom{000.0}\ \mathrm{ns}\ \ge\ \SI{450}{\nano\second}}.
\end{equation}

\subsection{LCD result}
The final code writes ASCII \texttt{0x41}; using the working LCD module, the character ``A'' was displayed clearly.

\figureplaceholder{figures/task5_LCD_A.jpg}
{Photograph the working LCD module clearly displaying the letter A.}
{Final LCD output displaying the character ``A''.}

% ============================================================
\clearpage
\section*{References}
\addcontentsline{toc}{section}{References}

\begin{enumerate}
  \item EEE3096S Practical 1B, \emph{Execution Timing and Cycle-Accurate Assembly on the UCT Dev Board}, UCT, 2026.
  \item STMicroelectronics, \emph{RM0091 STM32F0x1/STM32F0x2/STM32F0x8 Reference Manual}. \textbf{TO FILL:} exact clock-tree and timer sections used in Tasks 2--3.
  \item STMicroelectronics, \emph{STM32F051x4/x6/x8 Datasheet}. \textbf{TO FILL:} exact tables/sections used.
  \item Hitachi, \emph{HD44780U LCD Controller/Driver Datasheet}. Relevant items include the 4-bit initialisation flow and bus timing characteristics. \textbf{TO FILL:} exact figure/table/page references used in the final report.
  \item UCT Dev Board schematic and pinout reference. \textbf{TO FILL:} schematic sheet/reference used for the level-shifter resistor value and signal names.
  \item \textbf{TO FILL:} PC CPU documentation/source used for the Task 1 CPU clock, if cited.
\end{enumerate}

% ============================================================
\clearpage
\appendix
\section{Code Appendix}

The practical requires the report to be followed by a code appendix. Put the final source files in a local \texttt{code/} directory beside this \LaTeX{} file before compiling.

\subsection{Task 1 PC golden implementation}
\codefile{code/task1_golden.c}{Complete Task 1 golden-measure source.}

\subsection{Tasks 2--3 board C implementation}
\codefile{code/task2_3_main.c}{Complete board implementation and timing code used for Tasks 2 and 3.}

\subsection{Task 4 assembly}
\codefile{code/dsp.s}{Final Task 4 \texttt{dsp.s}.}

\subsection{Task 5 assembly}
\codefile{code/lcd.s}{Final Task 5 \texttt{lcd.s}.}

\section{Submission Evidence Checklist}
Before generating the final PDF, confirm all of the following:
\begin{itemize}
  \item The report has five numbered sections in the required order, followed by references and the code appendix.
  \item Task 2 has one full-screen scope capture of the PC13 single-call pulse.
  \item Task 3 has four full-screen captures: -O0, -O1, -O2 and -Os.
  \item Task 4 has one full-screen capture with CH1=PB0, CH2=PA4 and cursors on the measured phase delay.
  \item Task 5 includes before-fix and after-fix Enable captures; keep a separate rise-time capture as well if it is used for a separate numerical claim.
  \item Every scope image is uncropped and shows channel numbers, V/div, time/div, both cursors and the delta readout.
  \item Every timing value in the body text matches the corresponding oscilloscope readout.
  \item Every figure has a number, a caption naming the probed pins, the measured quantity and the measured value, and is referenced from the body text.
  \item The final LCD photograph clearly shows ``A'' on the working module.
  \item All exact manual sections, datasheet tables and external sources are included in the References section.
\end{itemize}

\end{document}
